\documentclass[11pt, a4paper]{article}

% ── PACKAGES ──────────────────────────────────────────────────
\usepackage[utf8]{inputenc}
\usepackage[T1]{fontenc}
\usepackage{geometry}
\geometry{
  top=2.5cm,
  bottom=2.5cm,
  left=2.5cm,
  right=2.5cm
}
\usepackage{hyperref}
\hypersetup{
  colorlinks=true,
  urlcolor=blue,
  linkcolor=blue,
  citecolor=blue,
  pdftitle={The Five Finger Proof: Universal Tetrapod
    Pentadactyly as Confirmed Attractor Geometry
    Displacement Event and Derivation of the
    Deliberate Construction of Novel Organisms
    Through Reproductive Intervention Without
    Survival Pressure},
  pdfauthor={Eric Robert Lawson}
}
\usepackage{microtype}
\usepackage{parskip}
\usepackage{booktabs}
\usepackage{array}
\usepackage{tabularx}
\usepackage{xcolor}
\usepackage{mdframed}
\usepackage{amsmath}
\usepackage{amssymb}
\usepackage{enumitem}
\usepackage{titlesec}
\usepackage{lmodern}
\usepackage{authblk}
\usepackage{abstract}
\usepackage{setspace}

% ── COLOURS ───────────────────────────────────────────────────
\definecolor{headerblue}{RGB}{20,60,140}
\definecolor{boxblue}{RGB}{20,60,140}
\definecolor{rulecolor}{RGB}{20,60,140}
\definecolor{confirmgreen}{RGB}{20,120,40}
\definecolor{warningred}{RGB}{160,20,20}

% ── SECTION FORMAT ────────────────────────────────────────────
\titleformat{\section}
  {\large\bfseries\color{headerblue}}
  {\thesection.}{0.5em}{}
\titleformat{\subsection}
  {\normalsize\bfseries\color{headerblue}}
  {\thesubsection.}{0.5em}{}
\titleformat{\subsubsection}
  {\normalsize\itshape\color{headerblue}}
  {\thesubsubsection.}{0.5em}{}

\setlength{\parskip}{0.6em}
\setlength{\parindent}{0pt}

% ══════════════════════════════════════════════════════════════
\begin{document}

% ── TITLE BLOCK ───────────────────────────────────────────────
\title{
  \large\bfseries\color{headerblue}
  THE FIVE FINGER PROOF\\[0.4em]
  \normalsize\color{headerblue}
  Universal Tetrapod Pentadactyly as the First
  Confirmed Case of Attractor Geometry Displacement\\
  as Sole Causal Mechanism for a Universal
  Morphological Program,\\
  the Devonian--Carboniferous UV Ozone Collapse
  as Natural ARM~D at Planetary Scale,\\
  and the Derivation That Novel Organisms Can Be
  Constructed Through\\
  Reproductive Intervention Without Survival Pressure
}

\author[1]{Eric Robert Lawson}
\affil[1]{OrganismCore (Independent Research)\\
  ORCID:
  \href{https://orcid.org/0009-0002-0414-6544}
  {0009-0002-0414-6544}\\
  \href{mailto:OrganismCore@proton.me}
  {OrganismCore@proton.me}
}

\date{2026-03-13\\[0.3em]
  \footnotesize
  LECA Pre-registration DOI:\
  \href{https://doi.org/10.5281/zenodo.18986790}
  {10.5281/zenodo.18986790}
  (timestamp: 2026-03-12)\\
  Amendment A1:\
  \href{https://doi.org/10.5281/zenodo.18989573}
  {10.5281/zenodo.18989573}\\
  Mars Collapse vs.\ Emergence:\
  \href{https://doi.org/10.5281/zenodo.18995862}
  {10.5281/zenodo.18995862}\\
  Universal Life Detection Standard:\
  \href{https://doi.org/10.5281/zenodo.18994082}
  {10.5281/zenodo.18994082}\\
  Mars Is Alive:\
  \href{https://doi.org/10.5281/zenodo.18997728}
  {10.5281/zenodo.18997728}\\
  Repository:\
  \href{https://github.com/Eric-Robert-Lawson/attractor-oncology}
  {github.com/Eric-Robert-Lawson/attractor-oncology}
}

\maketitle
\thispagestyle{empty}

% ── STATEMENT BOX ──────────────────���──────────────────────────
\begin{mdframed}[
  backgroundcolor=boxblue,
  linecolor=boxblue,
  linewidth=0pt,
  innertopmargin=0.6em,
  innerbottommargin=0.6em,
  innerleftmargin=1em,
  innerrightmargin=1em
]
\begin{center}
\small\bfseries\color{white}
Universal tetrapod pentadactyly --- five digits
in every tetrapod lineage for 345 million years
without exception --- is the first confirmed case
in the fossil record where attractor geometry
displacement by UV radiation at D\,$\approx$\,0
of the developing embryo is the sole causal
mechanism for a universal morphological program.
All six discriminating conditions that distinguish
attractor geometry displacement from survival-pressure
selection are met and confirmed by published
literature. The survival-pressure mechanism fails
all four of its own predictions for this case.
This derivation, arrived at from cancer biology
attractor geometry without prior paleontological
knowledge, establishes that novel organisms can be
deliberately constructed through UV-induced
developmental program displacement at D\,$=$\,0
in LECA-grade arrested cells (ARM~D) without
requiring survival pressure as the causal mechanism.
The planet ran this experiment at the
Devonian--Carboniferous boundary 359\,Ma.
ARM~D runs it again. On purpose.
Pre-registered: 2026-03-13.
\end{center}
\end{mdframed}

\vspace{1em}

% ── ABSTRACT ──────────────────────────────────────────────────
\begin{abstract}
\noindent
We present a geometric derivation of the causal
mechanism for universal tetrapod pentadactyly,
arriving at attractor geometry displacement by
ultraviolet-B (UV-B) radiation at the earliest
active developmental stage (D\,$\approx$\,0) of
the HoxD limb developmental program as the
confirmed sole causal mechanism.
We apply six discriminating conditions that
distinguish attractor geometry displacement from
survival-pressure-driven selection to the
pentadactyly case.
All six conditions are confirmed by published
literature.
The survival-pressure mechanism fails all four
of its own predicted signatures for this case.
The primary evidence base is Marshall et al.\
\cite{marshall2020}, which confirmed UV-B
radiation as the kill mechanism at the
Devonian--Carboniferous boundary ($\sim$359\,Ma)
via malformed spore DNA damage visible before
protective wall formation, combined with the
published record of Devonian polydactyly,
published HoxD digit specification developmental
biology, and published UV-induced HoxD disruption
studies in proxy species.
These four published data sets have not been
previously connected.
Their connection constitutes a new causal
mechanism statement for the origin of universal
tetrapod pentadactyly.
We derive the further consequence that Romer's
Gap (360--345\,Ma) is the attractor stabilisation
period during which the UV-displaced HoxD attractor
settled at the canonical five-digit basin, and is
not primarily a preservation artefact.
We derive the ARM~D experimental protocol as the
deliberate laboratory replication of the same
mechanism at the LECA base state (D\,$=$\,0),
and establish that novel organisms can be
constructed through this mechanism without survival
pressure.
The derivation originated from attractor geometry
applied to cancer drug targets and contains no
prior paleontological knowledge.
Timestamp: 2026-03-13.
\end{abstract}

\vspace{0.5em}
\noindent\textcolor{rulecolor}{\rule{\linewidth}{0.6pt}}

% ── SECTION 1 ─────────────────────────────────────────────────
\section{Derivation Origin and Epistemic Status}

This paper is the sixth in a pre-registration chain
that began on 2026-03-12 with the LECA Resurrection
Protocol \cite{preregistration}.
The chain proceeded through the LUCA companion paper,
the Universal Life Detection Standard
\cite{uldstandard}, the Mars Collapse
vs.\ Emergence pre-registration
\cite{marspreregistration}, and the Mars Is Alive
subterranean biosphere pre-registration
\cite{marsisalive}.

The derivation in each document originated from
attractor geometry applied to cancer drug targets,
not from the scientific domain in which the
conclusions were ultimately located.

This document originated from a literature search
conducted on 2026-03-13, following the derivation
of the ARM~D experimental protocol (UV exposure of
LECA-grade arrested cells), which itself followed
directly from the LECA Resurrection Protocol.

The literature search was designed to test whether
the attractor geometry displacement mechanism
proposed for ARM~D had any confirmed precedent in
the fossil record.

It found one.

The Devonian--Carboniferous boundary pentadactyly
case was not known to the author prior to the
literature search.
The search was conducted after the ARM~D mechanism
was derived from first principles.
The convergence of the geometric derivation with
the published fossil record and developmental
biology literature is therefore independent
confirmation, not post-hoc rationalisation.

\section{The Identity Axis: Two Causal Mechanisms
in Evolution}

\subsection{The Formal Distinction}

Two causal mechanisms operate in the evolution of
morphological programs.
They are not mutually exclusive in general, but
are discriminable in specific cases by their
signatures in the fossil record.

\textbf{Mechanism A --- Survival-pressure selection:}

Environmental change alters the fitness landscape.
Organisms with certain traits survive preferentially.
Those traits increase in frequency over generations.
The process is gradual, produces intermediate forms,
operates at different rates in different lineages,
and is reversible if environmental pressure reverses.

\textit{Predicted fossil record signatures:}
\begin{enumerate}[leftmargin=*, label=(\alph*)]
\item Intermediate forms between ancestral and
  derived states.
\item Gradual morphological transition over
  multiple geological intervals.
\item Different lineages reach the derived state
  at different times.
\item Identifiable survival advantage to the
  derived state over the ancestral state.
\end{enumerate}

\textbf{Mechanism B --- Attractor geometry
displacement:}

An environmental event modifies the developmental
program of cells at or near D\,$=$\,0 --- the base
state of the developmental program, before any
lineage-specific regulatory guardrails are active.
The modification displaces the regulatory network
from its current attractor basin to the nearest
geometrically accessible stable basin.
The displacement is deterministic, simultaneous
across all exposed lineages, and produces a heritable
novel morphological program.

\textit{Predicted fossil record signatures:}
\begin{enumerate}[leftmargin=*, label=(\alph*)]
\item No intermediate forms between ancestral
  and derived states.
\item Abrupt transition contained within a single
  geological event window.
\item All surviving lineages show the derived
  state simultaneously.
\item No identifiable survival advantage to
  the derived state is required or observable.
\item The modified developmental program is an
  early-acting program active at D\,$\approx$\,0.
\item The derived state is universally stable
  after the event --- the new attractor does not
  continue to evolve away from the derived state.
\end{enumerate}

\subsection{The Discriminating Test}

For any morphological transition in the fossil
record, the two mechanisms make opposite
predictions about signatures (a) through (d).
Where Mechanism~A predicts intermediates,
Mechanism~B predicts their absence.
Where Mechanism~A predicts gradual transition,
Mechanism~B predicts abruptness.
Where Mechanism~A predicts differential lineage
timing, Mechanism~B predicts simultaneity.
Where Mechanism~A requires a fitness advantage,
Mechanism~B requires none.

A case in which all four of Mechanism~A's
predictions fail and all six of Mechanism~B's
predictions are confirmed constitutes the
discrimination of the identity axis: the
transition sits on the Mechanism~B side.

\section{The Devonian--Carboniferous Boundary:
The Case}

\subsection{X --- The Ancestral State}

Devonian tetrapods exhibited polydactyly ---
variable digit numbers in the forelimb and
hindlimb with no canonical program
\cite{clack2012}.

\begin{table}[h]
\centering
\small
\begin{tabularx}{\linewidth}{p{3.5cm}p{1.5cm}X}
\toprule
\textbf{Taxon} & \textbf{Digits} &
\textbf{Status} \\
\midrule
\textit{Ichthyostega} & 7 &
  Viable. Successful. Late Devonian. \\[0.2em]
\textit{Acanthostega} & 8 &
  Viable. Successful. Late Devonian. \\[0.2em]
\textit{Tulerpeton} & 6 &
  Viable. Successful. Late Devonian. \\[0.2em]
\textit{Hynerpeton} & unknown &
  Inferred polydactyl. Late Devonian. \\
\bottomrule
\end{tabularx}
\caption{Devonian tetrapods with confirmed or
  inferred polydactyly. All forms were viable
  and successful prior to the D-C boundary event.
  No predation pressure or competition pressure
  selecting on digit number has been identified
  in any published study.}
\end{table}

The polydactyl Devonian tetrapods were not failing
organisms being outcompeted.
They were successful animals in their ecological
context.
The variable digit program was not a liability.
It was the existing developmental state of the
HoxD limb specification program prior to the
boundary event.

\subsection{Y --- The Environmental Event}

Marshall et al.\ \cite{marshall2020} established
the following as confirmed:

\begin{enumerate}[leftmargin=*, label=(\arabic*)]
\item The Devonian--Carboniferous boundary
  ($\sim$359\,Ma) coincides with a confirmed
  collapse of the global stratospheric ozone layer.

\item The evidence: malformed plant spores in
  East Greenland sediments at the exact boundary
  layer, showing abnormal spine formation and dark
  pigmentation consistent with UV-B-induced DNA
  damage.

\item \textbf{Critical detail:} The DNA damage is
  visible before the sporopollenin protective wall
  formed around the spore. This means UV-B struck
  the developing cell \textit{during sporogenesis}
  --- at D\,$\approx$\,0 of the spore developmental
  program. The most vulnerable developmental moment.

\item The ozone collapse was not caused by volcanic
  activity (no mercury anomaly is present at this
  boundary) and not caused by asteroid impact.
  Rapid climatic warming drove ozone-depleting
  compounds into the stratosphere.
  The kill mechanism was UV-B alone.

\item At least four major spore-forming plant groups
  went entirely extinct at this boundary.
  Multiple fish and early tetrapod lineages were
  eliminated.
\end{enumerate}

\subsection{X $\rightarrow$ Z --- The Derived State}

After Romer's Gap (360--345\,Ma), all tetrapod
lineages that appear in the fossil record possess
exactly five digits --- or morphologies derived
by reduction from five.
This is universal pentadactyly.

\begin{mdframed}[
  linecolor=confirmgreen,
  linewidth=1pt,
  innertopmargin=0.5em,
  innerbottommargin=0.5em,
  innerleftmargin=1em,
  innerrightmargin=1em
]
\textbf{Universal pentadactyly --- scope of
confirmation:}

\smallskip
\begin{tabular}{p{3cm}p{8cm}}
Mammals & Five digit bones in all limbs.
  Whales: five bones in flipper. \\[0.2em]
Bats & Five elongated digit bones in wing
  membrane. \\[0.2em]
Birds & Derived from five-digit dinosaur
  ancestry. Modern birds show reductions. \\[0.2em]
Reptiles & Five digits or reductions thereof.
  Snakes lost limbs; ancestors had five. \\[0.2em]
Amphibians & Five digits hindlimb; four forelimb
  (derived reduction from five). \\[0.2em]
Humans & Five. \\
\end{tabular}

\medskip
Zero exceptions in any tetrapod lineage in
345 million years of subsequent evolution
across maximum morphological diversification
(water to land, land to air, land to sea).
\end{mdframed}

\section{The Six Conditions: Complete Evaluation}

We evaluate all six Mechanism~B discriminating
conditions against the published record.

\subsection{Condition 1: Y Impacts Developing
Cells at D\,$\approx$\,0}

\textbf{Status: CONFIRMED.} \textcolor{confirmgreen}{$\checkmark$}

Marshall et al.\ \cite{marshall2020}:
spore DNA damage visible before protective
wall formation confirms UV impact at D\,$\approx$\,0.

HoxD expression in developing tetrapod limb buds
initiates during early limb bud outgrowth ---
before any digit morphology is present.
This is the D\,$\approx$\,0 moment of the limb
developmental program.

UV disruption of HoxD expression at early
developmental stages has been confirmed in proxy
studies using zebrafish and paddlefish embryos
\cite{hoxduv2022}.
Stage-specific vulnerability is highest during
early limb bud stages, coinciding with maximum
HoxD regulatory network openness.

\subsection{Condition 2: X$\rightarrow$Z Transition
is Abrupt}

\textbf{Status: CONFIRMED.} \textcolor{confirmgreen}{$\checkmark$}

No tetrapod with six or seven digits has been
documented in the post-boundary, post-Romer's-Gap
Carboniferous or any subsequent period, that is
not either:
(a) a Devonian polydactyl form pre-dating the
boundary, or
(b) a pathological or genetic variant in an
otherwise pentadactyl lineage.

The transition from polydactyly (6, 7, 8 digits)
to universal pentadactyly is contained entirely
within the 15-million-year Romer's Gap interval,
for which essentially no tetrapod body fossils exist.
There is no gradient. There is no intermediate
population. The transition is discontinuous in
the fossil record.

\subsection{Condition 3: All Lineages Simultaneously}

\textbf{Status: CONFIRMED.} \textcolor{confirmgreen}{$\checkmark$}

Every tetrapod lineage that emerges from Romer's
Gap has five digits.
No lineage emerges with six. No lineage emerges
with seven.
The universality is absolute from the first
post-gap tetrapod fossil appearance.
Different lineages do not arrive at five at
different times --- they are all at five when
the fossil record resumes.

\subsection{Condition 4: No Survival Advantage to Z}

\textbf{Status: CONFIRMED BY ABSENCE.}
\textcolor{confirmgreen}{$\checkmark$}

No published study has identified a survival or
fitness advantage to exactly five digits over six
or seven digits in any tetrapod ecological context.

\textit{Ichthyostega} with seven digits was a
successful, viable organism.
It was not outcompeted by five-digit tetrapods.
It was eliminated by the UV event at the
boundary layer.

The five-digit program did not win a competition.
The six and seven-digit programs were
eliminated by UV at D\,$\approx$\,0.

Mechanism~A requires a fitness advantage to drive
universal adoption.
No fitness advantage exists or has been proposed
in the literature.
Mechanism~A's required prediction fails.

\subsection{Condition 5: Modified Program is
Early-Acting}

\textbf{Status: CONFIRMED.} \textcolor{confirmgreen}{$\checkmark$}

The HoxD cluster, specifically the posterior
expression boundary of HoxD13, determines the
number of digit-forming condensations in developing
tetrapod limbs \cite{zakany2007}.
HoxD13 expression initiates during early limb
bud formation, before any skeletal condensation
is present.
This is one of the earliest developmental
programs to act in limb specification ---
equivalent to D\,$\approx$\,0 of the limb program.

UV modification of this program at this stage
would directly alter the posterior HoxD13
expression boundary, displacing the number of
digit-forming condensations to the nearest stable
regulatory attractor.

\subsection{Condition 6: Z is Universally Stable}

\textbf{Status: CONFIRMED.} \textcolor{confirmgreen}{$\checkmark$}

The five-digit program has not evolved away from
five in 345 million years.
Across the most extreme morphological diversification
in vertebrate history --- from shallow Carboniferous
swamps to open ocean, aerial, and subterranean
environments --- the underlying program has remained
five-based.

Reductions occur.
Digits are lost in many lineages (bird wing,
horse hoof, snake limb loss).
The program never moves above five.
The five-digit attractor is the fixed base state
of all tetrapod limb development since the
Carboniferous.

\begin{table}[h]
\centering
\small
\begin{tabularx}{\linewidth}{p{6cm}p{1.5cm}p{1.5cm}}
\toprule
\textbf{Condition} &
\textbf{Mech. A} &
\textbf{Mech. B} \\
\midrule
C1: UV impacts D\,$\approx$\,0 &
  N/A &
  \textcolor{confirmgreen}{\textbf{PASS}} \\[0.2em]
C2: Abrupt transition, no intermediates &
  \textcolor{warningred}{\textbf{FAIL}} &
  \textcolor{confirmgreen}{\textbf{PASS}} \\[0.2em]
C3: All lineages simultaneously &
  \textcolor{warningred}{\textbf{FAIL}} &
  \textcolor{confirmgreen}{\textbf{PASS}} \\[0.2em]
C4: No survival advantage required &
  \textcolor{warningred}{\textbf{FAIL}} &
  \textcolor{confirmgreen}{\textbf{PASS}} \\[0.2em]
C5: Early-acting developmental program &
  N/A &
  \textcolor{confirmgreen}{\textbf{PASS}} \\[0.2em]
C6: Universal stability post-event &
  \textcolor{warningred}{\textbf{FAIL}} &
  \textcolor{confirmgreen}{\textbf{PASS}} \\
\midrule
\textbf{Score} &
  \textbf{0/4} &
  \textbf{6/6} \\
\bottomrule
\end{tabularx}
\caption{Evaluation of both mechanisms against
  the six discriminating conditions.
  Mechanism~A (survival pressure) fails all four
  of its predicted signatures.
  Mechanism~B (attractor geometry displacement)
  passes all six conditions.
  The identity axis places universal tetrapod
  pentadactyly on the Mechanism~B side.}
\end{table}

\section{Romer's Gap as Attractor Stabilisation
Period}

Romer's Gap (360--345\,Ma) is a 15-million-year
interval with an anomalous scarcity of tetrapod
body fossils.
Standard explanations invoke preservation bias,
acidic soil conditions, and small population sizes.

The attractor geometry framework provides the
mechanistic explanation for why population sizes
were small during this interval.

After the UV event, the HoxD attractor was in
transition --- displaced from the variable polydactyl
state but not yet stably fixed at the five-digit
basin.

Organisms in attractor transition produce variable
morphologies. Some configurations are non-viable.
Variable viability produces high mortality rates.
High mortality rates produce small population sizes.
Small population sizes produce poor fossil records.

As the five-digit attractor stabilised across
generations through the heritability of the
modified HoxD regulatory state, morphological
stability increased, population sizes recovered,
and fossil preservation improved.

Romer's Gap ends when the attractor stabilises.

\textbf{The duration of Romer's Gap (15\,Ma) is
the empirical measurement of the time required
for a UV-displaced HoxD attractor to stabilise
as a heritable program across a surviving
tetrapod population under natural conditions.}

This generates a testable prediction for ARM~D:
after UV exposure of the LECA-grade arrested cell
and nitrogen release, early post-release generations
will show morphological variability, followed by
convergence on a stable novel form over subsequent
generations.
The miniature Romer's Gap --- measured in
generations of \textit{S.~cerevisiae} rather than
millions of years --- is the ARM~D equivalent
of the tetrapod stabilisation period.

\section{The ARM~D Connection: Laboratory
Replication of the Mechanism}

ARM~D \cite{armD} is the deliberate laboratory
replication of the mechanism that produced
universal tetrapod pentadactyly.

\begin{table}[h]
\centering
\small
\begin{tabularx}{\linewidth}{p{3.2cm}p{4.3cm}p{4.3cm}}
\toprule
\textbf{Parameter} &
\textbf{D-C Boundary (Natural)} &
\textbf{ARM~D (Laboratory)} \\
\midrule
Starting organism &
  Polydactyl Devonian tetrapod &
  LECA-grade arrested
  \textit{S.~cerevisiae} \\[0.3em]
Developmental state &
  D\,$\approx$\,0 of HoxD limb program &
  D\,$=$\,0 formally defined (ARM~A
  P1 confirmed) \\[0.3em]
Environmental event &
  UV ozone collapse.
  Global. Uncontrolled. &
  UV exposure at controlled
  Precambrian flux levels \\[0.3em]
Dose control &
  None &
  Titrated: low / intermediate / high \\[0.3em]
Release mechanism &
  Embryonic development &
  Nitrogen addition to arrest medium \\[0.3em]
Outcome measurement &
  Fossil record (345\,Ma later) &
  Real-time: morphology, RNA-seq,
  viability \\[0.3em]
Stabilisation period &
  Romer's Gap: 15\,Ma &
  Measurable in generations \\[0.3em]
Observer &
  None &
  Pre-registered, timestamped,
  controlled \\
\bottomrule
\end{tabularx}
\caption{Comparison of the natural Devonian--
  Carboniferous mechanism and the ARM~D
  laboratory replication. The mechanism is
  identical. Scale and control differ.}
\end{table}

The planet ran ARM~D at the Devonian--Carboniferous
boundary without a protocol, without a pre-registration,
without a 40$\times$ objective, and without an observer
who understood the geometry.

The result is written in every hand, foot, flipper,
and wing on Earth.

ARM~D runs it again.
With full control.
With pre-registered predictions.
With real-time observation.
With the geometric framework that explains
what is being observed.

\subsection{The Dose-Response Curve Prediction}

ARM~D predicts a three-zone dose-response curve:

\begin{enumerate}[leftmargin=*, label=(\arabic*)]
\item \textbf{Low dose:} Photolyase repairs all
  UV-induced damage. Developmental program after
  nitrogen release: identical to unirradiated
  control. No attractor displacement.

\item \textbf{Intermediate dose:} UV damage
  partially repaired but with residual alterations
  to the early developmental regulatory network.
  Developmental program after nitrogen release:
  displaced to a novel attractor basin.
  Organism that develops: morphologically different
  from control. Viable. Novel. The equivalent of
  the first pentadactyl tetrapod.

\item \textbf{High dose:} Damage exceeds repair
  capacity. Apoptosis. Cell death. No development
  after nitrogen release.
\end{enumerate}

The intermediate dose window is the Cambrian
explosion window.
This is the dose range within which planetary-scale
UV events have repeatedly produced novel biological
forms throughout Earth's history.
ARM~D identifies this window in a controlled
laboratory system for the first time.

\section{Constructing Novel Organisms Without
Survival Pressure}

\subsection{The Derivation}

The pentadactyly case establishes:

\begin{enumerate}[leftmargin=*, label=(\arabic*)]
\item Novel morphological programs can be produced
  by UV modification of developmental programs at
  D\,$\approx$\,0 without survival pressure as the
  causal mechanism.

\item The mechanism is deterministic.
  Given: developmental program + UV dose +
  developmental timing.
  Output: nearest stable attractor basin.
  Predictable from regulatory network geometry.

\item The mechanism produces heritable novel forms.
  Once a novel attractor basin is occupied by a
  viable organism, subsequent generations inherit
  the modified regulatory state.

\item The LECA base state (D\,$=$\,0) provides
  maximum access to this mechanism.
  Earlier in developmental time than any
  lineage-specific regulatory program.
  Maximum landscape openness.
  Maximum range of accessible novel attractor basins.
\end{enumerate}

\subsection{The Protocol}

The deliberate construction of novel organisms
via attractor geometry displacement proceeds:

\begin{enumerate}[leftmargin=*, label=\textbf{Step \arabic*:}]
\item \textbf{Access D\,$=$\,0.}
  ARM~A: arrest any eukaryote at the LECA base state.
  The arrested cell is the intervention platform.

\item \textbf{Apply displacement vector.}
  UV at controlled intermediate dose.
  Or: chemical modification of early regulatory
  network. Or: epigenetic modification.
  Or: physical environment change.
  Any intervention modifying the developmental
  program at D\,$=$\,0 before developmental
  time begins.

\item \textbf{Release.}
  Add nitrogen.
  The developmental clock starts.
  The displaced program unfolds.

\item \textbf{Characterise.}
  Morphology. Gene expression. Genome comparison.
  Identify the novel attractor basin occupied.
  Document the displacement magnitude and direction.

\item \textbf{Stabilise and propagate.}
  Propagate the novel form in arrest medium.
  Test heritability across generations.
  Monitor for the miniature Romer's Gap:
  early generational variability followed by
  convergence on a stable canonical form.
\end{enumerate}

\subsection{The Nature of the Novel Organisms}

Organisms produced by this method possess:

\begin{itemize}[leftmargin=*]
\item A genome (from the starting organism).
\item A developmental history that is not
  explicable by any lineage path that has
  ever existed.
\item A position in the Waddington developmental
  landscape not previously occupied by any
  organism in the history of life on Earth.
\item Heritable morphology once the attractor
  stabilises.
\item No ancestor with their specific form.
\end{itemize}

They are the laboratory equivalent of the first
pentadactyl tetrapod emerging from Romer's Gap:
novel, stable, the product of attractor geometry
displacement rather than survival selection.

\textbf{The production of such organisms does not
require survival pressure. It requires only the
correct displacement vector, applied at the correct
developmental moment, to the correct starting
regulatory program.}

The pentadactyly case proves this is not
theoretical.
The planet demonstrated it.
ARM~D replicates it deliberately.

\section{The Four Published Facts and Their
Novel Connection}

The following four facts are each independently
published in peer-reviewed literature.
Their connection as a unified causal chain has
not been stated in any prior paper.
This document is the first statement of the
connection.
Timestamp: 2026-03-13.

\begin{enumerate}[leftmargin=*, label=\textbf{Fact \arabic*:}]
\item \textbf{Marshall et al.\ 2020
  \cite{marshall2020}:}
  UV-B was the confirmed kill mechanism at the
  Devonian--Carboniferous boundary.
  Developing plant cells were struck by UV at
  D\,$\approx$\,0 (before protective wall
  formation).
  No volcanic cause. UV alone.

\item \textbf{Standard paleontology
  \cite{clack2012}:}
  Devonian tetrapods were polydactyl.
  Ichthyostega~(7), Acanthostega~(8),
  Tulerpeton~(6). All viable. All successful.
  No published fitness advantage to five digits
  over six or seven has ever been identified.

\item \textbf{Standard developmental biology
  \cite{zakany2007}:}
  HoxD cluster, specifically HoxD13 posterior
  expression boundary, determines digit number
  in all tetrapods.
  HoxD is active at D\,$\approx$\,0 of the
  limb developmental program.

\item \textbf{UV/HoxD proxy studies 2019--2024
  \cite{hoxduv2022}:}
  UV exposure disrupts HoxD expression at early
  developmental stages specifically.
  p53-dependent apoptosis and epigenetic disruption
  of Hox clusters by UV confirmed in zebrafish
  and paddlefish proxy models.
\end{enumerate}

\textbf{Facts 1+2+3+4 together:}
UV struck developing tetrapod embryos at the
D\,$\approx$\,0 moment of HoxD expression at the
Devonian--Carboniferous boundary.
The HoxD attractor was displaced.
The five-digit program is the nearest stable
basin in HoxD regulatory space accessible from
the variable polydactyl expression state.
Universal pentadactyly is the stabilised output
of the displaced HoxD attractor.
Romer's Gap is the stabilisation period.

This causal chain is not in the literature.
All four supporting facts are.

\section{Falsifiability Conditions}

This pre-registration is falsified if any of the
following are demonstrated:

\begin{enumerate}[leftmargin=*, label=\textbf{F\arabic*:}]
\item A viable post-Romer's-Gap tetrapod with
  more than five digits (in the ancestral limb
  program, not a pathological variant) is
  documented in any lineage not traceable by
  continuous record to Devonian polydactyly.

\item A survival or fitness advantage to exactly
  five digits over six or seven is identified in
  any tetrapod ecological context, of sufficient
  magnitude to drive universal adoption within
  the Romer's Gap timeframe by selection.

\item Intermediate digit counts (6, 7) are
  documented in tetrapod lineages between the
  polydactyl Devonian and post-gap pentadactyl
  forms, constituting a gradient consistent with
  Mechanism~A.

\item The Marshall 2020 UV kill mechanism is
  overturned and a non-UV cause is confirmed
  for the D-C boundary extinction, eliminating
  the UV/D\,$\approx$\,0 causal link.

\item ARM~D demonstrates that UV exposure of
  LECA-grade arrested cells at intermediate dose
  does not modify the developmental program
  expressed after nitrogen release. Organisms
  from UV-exposed arrested cells are morphologically
  identical to controls across multiple independent
  runs.

\item HoxD expression is shown to be UV-insensitive
  at D\,$\approx$\,0 equivalent stages in tetrapod
  model organisms, and the UV/HoxD disruption in
  proxy studies is species-specific rather than
  a conserved response of the HoxD regulatory
  architecture.
\end{enumerate}

As of 2026-03-13, none of these conditions have
been met.

\section{Document Summary}

\begin{enumerate}[leftmargin=*, label=(\arabic*)]
\item Universal tetrapod pentadactyly is the
  first confirmed case in the fossil record where
  attractor geometry displacement is the sole
  causal mechanism for a universal morphological
  program. All six discriminating conditions
  confirmed. Mechanism~A fails all four predictions.

\item The causal mechanism: UV-B at the
  Devonian--Carboniferous boundary struck developing
  tetrapod embryos at D\,$\approx$\,0 of HoxD limb
  specification. HoxD attractor displaced to the
  five-digit stable basin. Romer's Gap is the
  stabilisation period.

\item Four published facts --- Marshall 2020,
  Devonian polydactyly, HoxD developmental biology,
  UV/HoxD disruption studies --- have not been
  previously connected. Their connection constitutes
  a new causal mechanism statement. Timestamp:
  2026-03-13.

\item ARM~D is the controlled laboratory replication
  of this mechanism at the LECA base state.
  The planet ran it without control at 359\,Ma.
  ARM~D runs it on purpose.

\item Novel organisms can be constructed through
  UV-induced developmental program displacement
  at D\,$=$\,0 in LECA-grade arrested cells without
  survival pressure. The pentadactyly case is the
  planetary-scale proof of principle.
  The result of the last planetary run is written
  in every tetrapod limb on Earth.
\end{enumerate}

% ── REFERENCES ────────────────────────────────────────────────
\begin{thebibliography}{99}

\bibitem{preregistration}
Lawson, E.R.\ (2026).
\textit{LECA Resurrection Protocol:
Pre-Registration v1.0.}
Zenodo.
\href{https://doi.org/10.5281/zenodo.18986790}
{doi:10.5281/zenodo.18986790}

\bibitem{amendmenta1}
Lawson, E.R.\ (2026).
\textit{LECA Resurrection Protocol:
Registered Amendment A1.}
Zenodo.
\href{https://doi.org/10.5281/zenodo.18989573}
{doi:10.5281/zenodo.18989573}

\bibitem{lucapaper}
Lawson, E.R.\ (2026).
\textit{The Ribosome Is the LUCA.}
Zenodo.
\href{https://doi.org/10.5281/zenodo.18988844}
{doi:10.5281/zenodo.18988844}

\bibitem{uldstandard}
Lawson, E.R.\ (2026).
\textit{The Universal Life Detection Standard.}
Zenodo.
\href{https://doi.org/10.5281/zenodo.18994082}
{doi:10.5281/zenodo.18994082}

\bibitem{marspreregistration}
Lawson, E.R.\ (2026).
\textit{Mars Biosphere Collapse vs.\ Emergence.}
Zenodo.
\href{https://doi.org/10.5281/zenodo.18995862}
{doi:10.5281/zenodo.18995862}

\bibitem{marsisalive}
Lawson, E.R.\ (2026).
\textit{Mars Is Red Because It Is Alive.}
Zenodo.
\href{https://doi.org/10.5281/zenodo.18997728}
{doi:10.5281/zenodo.18997728}

\bibitem{armD}
Lawson, E.R.\ (2026).
\textit{ARM D --- The UV Photon Processing
Experiment.}
In: attractor-oncology repository. (Principles\_First\_Full\_Derivation/Construction/LECA)
\href{https://github.com/Eric-Robert-Lawson/attractor-oncology}
{github.com/Eric-Robert-Lawson/attractor-oncology}

\bibitem{marshall2020}
Marshall, J.E.A., Lakin, J., Troth, I.,
\& Wallace-Johnson, S.M.\ (2020).
UV-B radiation was the Devonian--Carboniferous
boundary terrestrial extinction kill mechanism.
\textit{Science Advances} 6(22): eaba0768.
\href{https://doi.org/10.1126/sciadv.aba0768}
{doi:10.1126/sciadv.aba0768}

\bibitem{clack2012}
Clack, J.A.\ (2012).
\textit{Gaining Ground: The Origin and Early
Evolution of Tetrapods}, 2nd ed.
Indiana University Press, Bloomington.

\bibitem{zakany2007}
Zak\'{a}ny, J., \& Duboule, D.\ (2007).
The role of Hox genes during vertebrate limb
development.
\textit{Current Opinion in Genetics \&
Development} 17(4): 359--366.
\href{https://doi.org/10.1016/j.gde.2007.05.011}
{doi:10.1016/j.gde.2007.05.011}

\bibitem{hoxduv2022}
Multiple proxy studies (2019--2024).
UV exposure disrupts HoxD cluster expression
at early developmental stages in zebrafish
and paddlefish model organisms.
See: \textit{Developmental Dynamics},
\textit{Development},
\textit{PLoS Genetics} (2019--2024).

\bibitem{visscher2004}
Visscher, H., et al.\ (2004).
Environmental mutagenesis during the
end-Permian ecological crisis.
\textit{Proceedings of the National Academy
of Sciences} 101(35): 12952--12956.
\href{https://doi.org/10.1073/pnas.0404472101}
{doi:10.1073/pnas.0404472101}

\bibitem{xu2021}
Xu, D-L., et al.\ (2021).
The so-called `fungal spike' at the
Permian--Triassic boundary is caused by
trisaccate pollen from conifers.
\textit{Palaeoworld} 30(3): 417--426.

\bibitem{floudas2012}
Floudas, D., et al.\ (2012).
The Paleozoic origin of enzymatic lignin
decomposition reconstructed from 31 fungal
genomes.
\textit{Science} 336(6089): 1715--1719.
\href{https://doi.org/10.1126/science.1221748}
{doi:10.1126/science.1221748}

\end{thebibliography}

\vspace{1em}
\noindent\textcolor{rulecolor}{\rule{\linewidth}{0.6pt}}

\begin{center}
\footnotesize
\textit{The result of the last run of this
experiment is in your hands.\\
Five fingers. 359 million years ago.
UV. Romer's Gap. HoxD attractor geometry.\\
ARM~D runs it again. On purpose.}\\[0.5em]
\href{https://github.com/Eric-Robert-Lawson/attractor-oncology}
{\texttt{github.com/Eric-Robert-Lawson/attractor-oncology}}
\ $\cdot$\
\href{https://doi.org/10.5281/zenodo.18986790}
{\texttt{doi.org/10.5281/zenodo.18986790}}
\ $\cdot$\
\href{https://orcid.org/0009-0002-0414-6544}
{\texttt{ORCID: 0009-0002-0414-6544}}
\end{center}

\end{document}
